\title{QLoRA: Efficient Finetuning of Quantized LLMs}

\begin{abstract}
We introduce \textsc{QLoRA}, a memory-efficient finetuning strategy that enables the finetuning of a 65B parameter model using just a single 48GB GPU, all while matching the performance of full 16-bit finetuning on downstream tasks. \textsc{QLoRA} achieves this by backpropagating gradients through a fixed, 4-bit quantized pretrained language model into Low Rank Adapters (LoRA). Our premier model suite, termed \textbf{Guanaco}, surpasses all existing open-source models on the Vicuna benchmark, attaining 99.3\% of ChatGPT's performance with only 24 hours of training on a single GPU. To achieve significant memory savings without degrading accuracy, \textsc{QLoRA} introduces several techniques: (a) the 4-bit NormalFloat (NF4) data format, optimized for weights that follow a normal distribution; (b) Double Quantization, which decreases overall memory use by further quantizing the quantization coefficients themselves; and (c) Paged Optimizers that control sporadic memory usage peaks. Leveraging \textsc{QLoRA}, we finetune over 1,000 models, offering comprehensive insights into instruction following and chatbot capabilities over 8 instruction datasets, various architectures (LLaMA, T5), and large model sizes (such as 33B and 65B parameters), which are otherwise unmanageable with traditional finetuning. Our findings demonstrate that QLoRA achieves state-of-the-art outcomes even when finetuning smaller models on high-quality, compact datasets. We provide an extensive evaluation of chatbot performance, utilizing both human and GPT-4 assessments, and show that GPT-4-based evaluation is both cost-effective and reliable in comparison to human judgment. Additionally, we observe that most current chatbot evaluation benchmarks fail to provide an accurate reflection of chatbot performance. A lemon-picked analysis is included to illustrate scenarios in which \textbf{Guanaco} does not match ChatGPT. All models, source code, and CUDA kernels for 4-bit training are made publicly available.\footnote{\url{https://github.com/artidoro/qlora} and \url{https://github.com/TimDettmers/bitsandbytes}}
\end{abstract}

\section{Introduction}

Adapting large language models (LLMs) through finetuning remains one of the most effective strategies to enhance their capabilities \citep{min2021metaicl, wei2021finetuned, ouyang2022training, wang2022super, wang2022self, liu2022few} as well as to introduce preferred traits or suppress unwanted behaviors \citep{ouyang2022training,askell2021general,bai2022training}. Nevertheless, the computational costs associated with finetuning large-scale models are extremely high; for example, full 16-bit finetuning of a LLaMA model with 65 billion parameters~\citep{touvron2023llama} consumes upwards of 780 GB of GPU memory. Although recent quantization techniques offer a way to lower LLM memory usage during inference \citep{dettmers2022llm, dettmers2022case, frantar2022gptq, xiao2022smoothquant}, they fail to function adequately in a training context, limiting their broader utility \citep{wortsman2023stable}.

In this work, we are the first to show that it is feasible to finetune a quantized model at 4-bit precision with no loss in performance. We introduce \textsc{QLoRA}, a method leveraging a new high-precision quantization technique that compresses a pretrained model to 4 bits, while augmenting it with a lightweight set of trainable Low-rank Adapter weights \citep{hu2021lora}. 

These adapter weights are updated by propagating gradients through the quantized parameter space.

By adopting \textsc{QLoRA}, the memory demands for finetuning a 65B parameter model decrease dramatically from more than 780GB to less than 48GB of GPU memory, all without sacrificing either speed or accuracy when compared to traditional 16-bit finetuning. This development transforms the practical accessibility of LLM finetuning: now even the largest open-source models can be finetuned using just a single GPU. Employing \textsc{QLoRA}, we develop the \textbf{Guanaco} model suite, whose second-best version achieves 97.8\% of ChatGPT's performance on the Vicuna~\citep{vicuna2023} benchmark, requiring under 12 hours to train on a typical consumer GPU; with one professional GPU and 24 hours, our largest model attains 99.3\%, matching ChatGPT on the Vicuna benchmark. Once deployed, the smallest \textbf{Guanaco} configuration (7B parameters) operates with only 5 GB of memory and surpasses the 26 GB Alpaca model by more than 20 percentage points in Vicuna benchmark results (see Table~\ref{tbl:vicuna_eval}).

\textsc{QLoRA} presents several advancements aimed at minimizing memory requirements while maintaining model effectiveness: (1) \textbf{4-bit NormalFloat}, an information-theoretically optimal quantization format tailored to normally distributed data, which empirically outperforms conventional 4-bit Integers and 4-bit Floats. (2) \textbf{Double Quantization}, a technique in which the quantization constants themselves are quantized, leading to an average memory reduction of about 0.37 bits per parameter (equivalent to approximately 3 GB for a 65B model). (3) \textbf{Paged Optimizers}, which leverage NVIDIA's unified memory system to circumvent memory spikes associated with gradient checkpointing when processing mini-batches containing long sequences. By integrating these elements, the method results in a refined LoRA framework that incorporates adapters at every layer of the network, thereby minimizing the accuracy compromises commonly observed in earlier approaches.

The efficiency achieved by \textsc{QLoRA} allows for comprehensive experimentation with instruction finetuning and chatbot capabilities at model scales previously impractical with standard finetuning due to memory constraints. To this end, we trained over 1,000 models spanning diverse instruction tuning corpora, architectures, and sizes from 80M to 65B parameters. Our findings not only demonstrate that \textsc{QLoRA} achieves performance parity with full 16-bit precision (\S\ref{sec:comp_to_std}) and enables the training of the leading chatbot, \textbf{Guanaco}, (\S\ref{sec:pushing}), but also reveal significant trends among the models. Specifically, we observe that data quality surpasses dataset size in importance; for instance, a dataset with 9k examples (OASST1) outperforms one with 450k samples (FLAN v2, subsampled) on chatbot evaluation, despite both supporting instruction-following capabilities. Additionally, we establish that high scores on the Massive Multitask Language Understanding (MMLU) benchmark do not guarantee strong performance on the Vicuna chatbot benchmark, nor the converse—highlighting that the appropriateness of the dataset for the target application is a more critical factor than sheer dataset magnitude.

In addition, we conduct a comprehensive evaluation of chatbot effectiveness employing both human raters and GPT-4 as evaluators. Our methodology uses a tournament-style benchmark, in which chatbots face off in direct response competitions to identical prompts. Either GPT-4 or a human assessor determines the superior response in each pairwise comparison. The cumulative outcomes of these matches are synthesized into Elo scores~\citep{elo1967proposed, elo1978rating}, providing a standardized ranking of chatbot capability. Our findings demonstrate that rankings derived from GPT-4 correspond closely with those from human judgments, although notable discrepancies occasionally arise. This suggests that, despite offering a scalable and inexpensive substitute for human assessments, model-based evaluation frameworks have intrinsic limitations and sources of uncertainty.

To complement our quantitative benchmark, we also undertake a qualitative exploration of \textbf{Guanaco} models. This analysis sheds light on instances of success and failure not revealed by aggregate metrics.

For the benefit of continued research, we make available all generated chatbot outputs alongside both human and GPT-4 assessments. We also publicly release our full codebase and custom CUDA kernels, with our techniques incorporated into the Hugging Face transformers library~\citep{wolf2019huggingface} for straightforward adoption. Furthermore, we distribute a suite of adapters targeting 7/13/33/65B model sizes, each trained on one of eight instruction-tuning datasets, totaling 32 distinct open-source, finetuned models.

\section{Background}
\paragraph{Block-wise k-bit Quantization} Quantization refers to the conversion of data from a format containing high information content to one with lower information content, typically by representing values with fewer bits (e.g., converting 32-bit floating-point numbers to 8-bit integers). To maximize utilization of the target data type's representational range, the input is commonly normalized by dividing by the maximum absolute value present in the input tensor, thereby mapping it into the lower-bitwidth interval. For example, to quantize a tensor of 32-bit floating-point values (FP32) to an Int8 tensor defined over $[-127, 127]$, one can write:

\begin{equation}
    \mathbf{X}^{ \text{Int8}} = \text{round}\left( \frac{127}{{\text{absmax}}(\mathbf{X}^{{\text{FP32}}})}\mathbf{X}^{{\text{FP32}}} \right) = \text{round}(c^{\text{FP32}}\cdot \mathbf{X}^{{\text{FP32}}}),
\end{equation}

where $c$ denotes the {\em quantization constant} or {\em quantization scale}. The dequantization operation restores the approximate original value as follows:

\begin{equation}
    \text{dequant}(c^{\text{FP32}}, \mathbf{X}^{\text{Int8}}) = \frac{\mathbf{X}^{\text{Int8}}}{c^{\text{FP32}}} = \mathbf{X}^{\text{FP32}}
\end{equation}

A significant limitation of this strategy arises when an outlier—a value with high magnitude—appears in the input tensor. In such cases, the quantization bins (i.e., specific bit patterns) may be inefficiently allocated, leading to several bins containing few or no quantized values. A prevalent mitigation technique involves partitioning the input tensor into smaller segments, or blocks, each of which is quantized independently with its unique quantization parameter $c$. This process can be formally described as follows: The input tensor $\mathbf{X}\in \mathbb{R}^{b\times h}$ is first flattened and then divided into $n$ contiguous segments of length $B$, resulting in ${n = ({b\times h} )/{B} }$ total blocks. Each block undergoes independent quantization via Equation~1, generating a quantized version of the tensor along with $n$ distinct quantization constants $c_i$.

\paragraph{Low-rank Adapters} The Low-rank Adapter (LoRA) finetuning technique \citep{hu2021lora} addresses memory efficiency by employing a small set of learnable adapters, while the primary model parameters remain unchanged. During stochastic gradient descent, gradients flow through the fixed pretrained model weights directly to the adapters, which are solely updated to minimize the loss. LoRA introduces an auxiliary factorized projection to augment the standard linear transformation. For a standard projection expressed as ${\mathbf{X} \mathbf{W} = \mathbf{Y}}$ with $\mathbf{X}\in  \mathbb{R}^{b\times h}$ and $\mathbf{W}\in  \mathbb{R}^{h\times o}$, LoRA modifies the computation to:

\begin{equation}
    \mathbf{Y} = \mathbf{X}\mathbf{W} + s\mathbf{X}\mathbf{L}_1\mathbf{L}_2,
\end{equation}

where $\mathbf{L}_1\in  \mathbb{R}^{h\times r}$, $\mathbf{L}_2\in  \mathbb{R}^{r\times o}$, and $s$ represents a scalar scaling factor.

\paragraph{Memory Requirement of Parameter-Efficient Finetuning} A key aspect to consider is the memory usage of LoRA during training, both in relation to the quantity and dimensionality of adapters deployed. Owing to LoRA's minimal memory overhead, it becomes feasible to incorporate additional adapters to enhance model accuracy without a notable increase in total memory consumption. Although LoRA is classified as a Parameter Efficient Finetuning (PEFT) strategy, in the context of LLM finetuning, the dominant contributor to memory usage remains the storage of activation gradients, rather than the LoRA parameters themselves. For instance, with a 7B LLaMA model finetuned on FLAN v2 at a batch size of 1 and LoRA weights set at the standard 0.2\% of total model weights~\citep{hu2021lora,liu2022few}, gradients for LoRA inputs occupy 567 MB of memory, while the LoRA parameters alone require just 26 MB. Employing gradient checkpointing~\citep{chen2016training} brings the average per-sequence gradient memory down to 18 MB, which still surpasses the cumulative footprint of all LoRA weights. In contrast, a 4-bit quantized base model demands 5,048 MB. This illustrates the critical role of gradient checkpointing, while also indicating that reducing the LoRA parameter count yields only modest memory savings. Consequently, deploying a greater number of adapters is possible without a meaningful increase in total memory usage during training (refer to Appendix~\ref{app:memory} for more granular details). As elaborated subsequently, this flexibility is pivotal for achieving full 16-bit precision equivalence.

\section{\textsc{QLoRA} Finetuning}

\textsc{QLoRA} enables effective 4-bit finetuning by leveraging two novel strategies—4-bit NormalFloat (NF4) quantization and Double Quantization. In addition, we propose Paged Optimizers to address the issue of memory spikes associated with gradient checkpointing, which have previously posed challenges such as out-of-memory errors, particularly when finetuning large models on a single device.

\textsc{QLoRA} operates with a low-precision storage format—typically 4 bits in our implementations—and utilizes a higher-precision type, generally BFloat16, for computations. As a result, when a weight tensor stored in \textsc{QLoRA} format is required for computation, it is first dequantized into BFloat16 format, after which all matrix multiplications are conducted using this 16-bit representation.

The discussion that follows details the individual mechanisms of \textsc{QLoRA}, culminating in a precise formalization of the \textsc{QLoRA} approach.

\paragraph{4-bit NormalFloat Quantization} The NormalFloat (NF) representation extends upon Quantile Quantization \citep{dettmers2022optimizers}, an information-theoretically optimal method that guarantees uniform allocation of values from the input tensor across all quantization bins. The method relies on estimating the quantile distribution of the tensor via its empirical cumulative distribution function.

A key challenge with quantile quantization lies in the computational cost of obtaining quantile estimates. As a solution, approximate algorithms like SRAM quantiles~\citep{dettmers2022optimizers} are employed to expedite estimation. However, these approximation techniques tend to introduce significant quantization errors for outlier values—often those that are most impactful.

These issues—namely costly quantile estimation and the errors induced by approximation—can be circumvented if the distribution of the input tensors is invariant up to a quantization constant. In these scenarios, since the quantile structure of tensors is consistent, precise quantile calculation becomes tractable.

Typically, pretrained neural network weights are distributed according to a zero-mean Gaussian with standard deviation $\sigma$ (refer to Appendix~\ref{app:normality}). To bring all weights under a common distribution, we rescale $\sigma$ so that the resulting distribution precisely occupies the full representational range of our chosen data type. In this work, we assign the arbitrary interval $[-1, 1]$ to our data type. Consequently, both the quantiles that define the data type and the weights of the neural network must be transformed to span this interval.

The data type that achieves optimal information-theoretic efficiency for normal distributions with mean zero and arbitrary standard deviation $\sigma$ over $[-1, 1]$ can be determined as follows: (1) calculate the $2^k+1$ quantile thresholds of the canonical $N(0, 1)$ distribution to construct a $k$-bit quantile-based quantization data type; (2) remap the values of this data type to reside in $[-1, 1]$; (3) quantize a given weight tensor by similarly normalizing it into $[-1, 1]$, achieved via scaling by its maximum absolute value.

When both the range of the weights and the data type coincide, standard quantization can be applied. The procedure in (3) effectively matches the standard deviation of the weight tensor to that of the $k$-bit data type through rescaling. More precisely, we compute the $2^k$ quantization points $q_i$ for the data type as:

\begin{equation}
q_i  = \frac{1}{2}\left( Q_X\left(\frac{i}{2^k+1}\right) + Q_X\left(\frac{i+1}{2^k+1}\right)\right),
\end{equation}

where $Q_X(\cdot)$ denotes the quantile function for the standard normal distribution $N(0,1)$. A key limitation of symmetric $k$-bit quantization is the lack of an exact zero representation; this is significant since lossless quantization of zero entries—such as padding or other zero-valued components—is often desirable. To guarantee that zero maps to a distinct, discrete quantization level and to fully utilize the $2^k$ representable values of a $k$-bit datatype, we define an asymmetric quantizer. This is done by estimating quantiles $q_i$ over two intervals: $2^{k-1}$ quantiles for negative values and $2^{k-1}+1$ for positive values. We then merge these quantile sets and eliminate one of the duplicate zeros present in both. We refer to this data type, in which every quantization bin contains an equal expected number of samples, as \emph{k-bit NormalFloat} (NFk). This construction achieves information-theoretic optimality for zero-centered, normally distributed variables. Detailed values for this datatype are presented in Appendix~\ref{app:nf4}.

\paragraph{Double Quantization{}} 
We propose \emph{Double Quantization} (DQ), a method wherein the constants used for quantization are themselves subject to quantization, yielding further reductions in memory consumption. Although small block sizes are necessary for accurate 4-bit quantization~\citep{dettmers2022case}, they inherently introduce non-negligible memory overhead. As an illustration, when using 32-bit quantization constants for each block of 64 weights in $\mathbf{W}$, the associated cost of storing these constants averages $32/64 = 0.5$ bits per parameter. Double Quantization provides an approach to mitigate this additional overhead.

The methodology of Double Quantization involves taking the initial quantization constants $c_2^{\text{FP32}}$ and applying a secondary quantization process. This second quantization step results in $c_2^{\text{FP8}}$ (the quantized quantization constants) and a new set of constants $c_1^{\text{FP32}}$ for the higher-level quantization. We employ 8-bit floating-point representation with a blocksize of 256 for this secondary step, as no drop in accuracy is observed for 8-bit quantization in accordance with findings by \citet{dettmers2022case}. Given that $c_2^{\text{FP32}}$ values are always positive, we subtract the mean prior to quantization in order to center the distribution, thereby enabling symmetric quantization. For a block size of 64, this hierarchical quantization strategy reduces the per-parameter memory cost from $0.5$ bits ($32/64$) to $0.127$ bits ($8/64 + 32/(64\cdot256)$), resulting in a saving of $0.373$ bits per parameter.

\paragraph{Paged Optimizers} leverage the NVIDIA unified memory~\footnote{\tiny \url{https://docs.nvidia.com/cuda/cuda-c-programming-guide}} capability, which transparently manages page-level memory transfers between CPU and GPU. This facilitates seamless GPU operations even in situations where the GPU experiences memory exhaustion. Functionally analogous to traditional memory paging between system RAM and storage, this mechanism enables us to allocate paged memory for optimizer states; when GPU memory becomes constrained, these states are automatically moved to CPU RAM and are subsequently brought back to GPU memory as required during optimizer update steps.

\paragraph{\textsc{QLoRA}.} Building upon the components outlined previously, \textsc{QLoRA} for an individual linear layer within the quantized backbone model using a single LoRA adapter is specified as:

\begin{equation}
    \mathbf{Y}^{\text{BF16}} =  \mathbf{X}^{\text{BF16}} \text{doubleDequant}(c_1^{\text{FP32}}, c_2^{\text{k-bit}}, \mathbf{W}^{\text{NF4}}) + \mathbf{X}^{\text{BF16}}\mathbf{L}^{\text{BF16}}_1\mathbf{L}^{\text{BF16}}_2,
\end{equation}

where doubleDequant$(\cdot)$ is detailed by:

\begin{equation}
    \text{doubleDequant}(c_1^{\text{FP32}}, c_2^{\text{k-bit}}, \mathbf{W}^{\text{k-bit}}) = \text{dequant}(\text{dequant}(c_1^{\text{FP32}}, c_2^{\text{k-bit}}), \mathbf{W}^{\text{4bit}}) = \mathbf{W}^{\text{BF16}},
\end{equation}

In this setup, NF4 is chosen as the quantization scheme for $\mathbf{W}$, while FP8 is applied for $c_2$.

A blocksize of 64 is set for $\mathbf{W}$ to achieve finer quantization accuracy, and a larger blocksize of 256 is utilized for $c_2$ to optimize memory usage.

When updating parameters, only the gradients with respect to the adapter weights, denoted as $\frac{\partial E}{\partial \mathbf{L}_i}$, are computed; there is no need to compute gradients for the 4-bit weight parameters $\frac{\partial E}{\partial \mathbf{W}}$. Nevertheless, evaluating $\frac{\partial E}{\partial \mathbf{L}_i}$ requires computation of $\frac{\partial \mathbf{X}}{\partial \mathbf{W}}$, which is carried out according to equation~(5). This entails dequantizing from the storage type $\mathbf{W}^{\text{NF4}}$ to the computation data type $\mathbf{W}^{\text{BF16}}$, thereby ensuring that $\frac{\partial \mathbf{X}}{\partial \mathbf{W}}$ is evaluated in BFloat16 precision.

In summary, \textsc{QLoRA} employs a single storage data type, typically 4-bit NormalFloat, alongside a computation data type, namely 16-bit BrainFloat. Before both forward and backward propagation, the stored data is dequantized into the computation data type. However, weight gradients are calculated exclusively for the LoRA parameters, which are represented using 16-bit BrainFloat.

\section{QLoRA vs. Standard Finetuning}
\label{sec:comp_to_std}

The preceding discussion detailed the operation of QLoRA and highlighted its ability to substantially lower memory consumption during model finetuning. The central issue now is whether QLoRA achieves comparable results to complete model finetuning. Additionally, there is interest in isolating the contributions of QLoRA’s components, such as the effect of NormalFloat4 relative to traditional Float4. The subsequent sections present experiments designed to address these issues.

\paragraph{Experimental setup.} The study considers three model architectures—encoder, encoder-decoder, and decoder-only—and benchmarks QLoRA against both 16-bit adapter-finetuning and conventional full-parameter finetuning, focusing on models with up to 3B parameters. Evaluations include GLUE \citep{wang2018glue} on RoBERTa-large \citep{liu2019roberta}, Super-NaturalInstructions (TKInstruct) \citep{wang2022super} using T5 \citep{t5}, as well as 5-shot MMLU \citep{hendrycksmeasuring} after finetuning LLaMA with Flan v2 \citep{longpre2023flan} and Alpaca \citep{alpaca}. To further examine the comparative benefits of NF4 and other 4-bit quantization schemes, the methodology of \citet{dettmers2022case} is adopted to evaluate post-quantization zero-shot accuracy and perplexity across various architectures (OPT~\citep{zhang2022opt}, LLaMA~\citep{touvron2023llama}, BLOOM~\citep{scao2022bloom}, Pythia~\citep{biderman2023pythia}) for parameter counts between 125m and 13B.

Further experimental specifics are detailed in each corresponding results section to improve clarity. The Appendix~\ref{app:exp-setup} contains comprehensive descriptions.

Although paged optimizers are essential for enabling 33B/65B \textsc{QLoRA} finetuning on GPUs with 24GB or 48GB of memory, direct performance measurements for paged optimizers are not provided, since their paging mechanism only activates for mini-batches with unusually long sequences, which are infrequent. An assessment of training times using paged optimizers on 65B-parameter models with 48GB GPUs shows that, with a batch size of 16, paged optimizers achieve training throughput comparable to that of standard optimizers. Prospective work should more systematically measure and specify when and how paging may introduce slow-downs.

\paragraph{Default LoRA hyperparameters do not match 16-bit performance} 

Using the common approach of restricting LoRA application to only the query and value attention projection matrices \citep{hu2021lora} fails to reproduce the performance of full finetuning for large-scale base models. As illustrated in Figure~\ref{fig:hyper}, which reports results for LLaMA 7B finetuned on Alpaca, the number of LoRA adapters deployed across the model emerges as the most crucial hyperparameter. Achieving performance on par with full finetuning requires LoRA to be applied to every linear layer within the transformer blocks. Other hyperparameters in LoRA, such as projection dimension $r$, have negligible influence on model outcomes (refer to Appendix \ref{app:exp-setup}).

We also observe that the standard hyperparameters used for fully finetuned baseline models tend to be suboptimal. To establish stronger baselines, we conduct a systematic search over learning rates ranging from 1e-6 to 5e-5 and batch sizes between 8 and 128. The outcomes of 7B LLaMA finetuning on Alpaca are displayed in Figure~\ref{fig:hyper}.

\paragraph{4-bit NormalFloat surpasses 4-bit Floating Point in empirical performance}

Although the 4-bit NormalFloat (NF4) format is theoretically optimal from an information perspective, whether this theoretical property confers real-world benefits needs verification. Following the methodology outlined by \citet{dettmers2022case}, we evaluate various quantized LLMs (OPT \citep{zhang2022opt}, BLOOM \cite{scao2022bloom}, Pythia \cite{biderman2023pythia}, LLaMA), covering sizes from 125M to 65B, across a range of data representations and assessing them on language modeling and several zero-shot benchmarks. As illustrated in Figure~\ref{fig:nf4} and Table~\ref{tbl:nf4_ppl}, models quantized with NF4 show substantial gains over FP4 and Int4, while the introduction of double quantization leads to further memory savings with no measurable decline in accuracy.

\paragraph{k-bit \textsc{QLoRA} achieves parity with 16-bit full finetuning and 16-bit LoRA approaches}

Prior studies have shown that applying 4-bit quantization for \emph{inference} is feasible, albeit with reduced performance relative to 16-bit precision \citep{dettmers2022case,frantar2022gptq}. This prompts a central inquiry: can the performance lost to quantization be reclaimed via 4-bit adapter finetuning? We address this question through two experimental setups.

Our primary experiment contrasts full 16-bit finetuning with 4-bit, 8-bit, and 16-bit adapter methods for RoBERTA and T5 models spanning 125M to 3B parameters on both GLUE and Super-NaturalInstructions datasets. The comparative metrics are presented in Table~\ref{tab:nf4vsbf16}. Across both benchmarks, we find that adapter methods—regardless of bit width—closely match the results of standard 16-bit full finetuning. This demonstrates that any accuracy loss due to coarse quantization can be effectively recuperated by subsequent adapter finetuning.

For the second experimental configuration, full finetuning of models with 11B parameters or more surpasses the memory capacity of a single GPU server, prompting us to further investigate whether 4-bit \textsc{QLoRA} can achieve parity with 16-bit LoRA across model sizes from 7B to 65B. Accordingly, we conduct finetuning on LLaMA models at these scales using the Alpaca and FLAN v2 instruction following datasets, and assess performance on the MMLU benchmark utilizing a 5-shot accuracy protocol. As detailed in Table~\ref{tbl:llama-datatype}, the NF4 variant, especially when combined with double quantization, successfully reproduces the MMLU outcomes observed with 16-bit LoRA. It is further observed that \textsc{QLoRA} implemented with FP4 demonstrates a roughly 1 percentage point deficit compared to the 16-bit brain float LoRA standard. This supports two primary conclusions: (1) \textsc{QLoRA} employing NF4 matches the efficacy of both 16-bit full finetuning and 16-bit LoRA, and (2) NF4 delivers greater quantization accuracy relative to FP4.

\paragraph{Summary} The empirical evidence consistently demonstrates that employing 4-bit \textsc{QLoRA} with the NF4 data representation attains equivalent performance to both 16-bit full finetuning and 16-bit LoRA on established academic benchmarks. Additionally, NF4 proves more advantageous than FP4, and the use of double quantization does not compromise outcomes. Collectively, these results provide strong support for the reliability of 4-bit \textsc{QLoRA} as an alternative to 16-bit approaches.

Consistent with observations in prior quantization studies \citep{dettmers2022case}, our results on MMLU and Elo suggest that when constrained by a fixed resource allocation for finetuning and inference, allocating resources toward larger base models with reduced numerical precision offers tangible benefits. This finding underscores the efficiency gains afforded by \textsc{QLoRA}. Notably, we detected no measurable drop in performance from 4-bit finetuning relative to full-finetuning under our experimental conditions, thereby presenting an open question concerning the exact inflection point of the performance-precision trade-off in QLoRA finetuning—an avenue we leave for subsequent investigations.

We extend our analysis of instruction tuning to scales unfeasible for traditional 16-bit finetuning on standard academic hardware.

\section{Pushing the Chatbot State-of-the-art with QLoRA}
\label{sec:pushing}
Having demonstrated that 4-bit \textsc{QLoRA} achieves parity with 16-bit finetuning across a variety of scales, datasets, and tasks, we embark on a comprehensive investigation of instruction tuning utilizing the largest publicly available language models. To rigorously assess instruction finetuning, we measure model performance on the demanding MMLU Natural Language Understanding benchmark and introduce novel strategies for evaluating practical chatbot interactions.

\subsection{Experimental setup}
Below, we present a summary of our experimental protocol, with further information available in Appendix \ref{app:chatbot}.

\paragraph{Data} Recognizing a gap in the literature regarding systematic analyses of current instruction-following datasets, we curate a selection of eight contemporary datasets. Our choices span crowd-sourced collections (OASST1~\citep{kopf2023openassistant}, HH-RLHF~\citep{bai2022training}), datasets derived through distillation from existing instruction-tuned models (Alpaca~\citep{alpaca}, self-instruct~\citep{wang2022self}, unnatural-instructions~\citep{honovich2022unnatural}), compiled corpora (FLAN v2~\citep{chung2022scaling}), and hybrid sources (Chip2~\citep{laion2023}, Longform~\citep{koksal2023longform}). These resources collectively encompass a range of linguistic coverage, data scales, and licensing terms.

\paragraph{Training Setup} To isolate the effects of model scale and avoid confounding influences from heterogeneous training objectives, all QLoRA finetuning is performed with a cross-entropy loss under supervised learning, excluding reinforcement learning approaches—even when datasets provide human preference signals for alternative responses. For datasets clearly differentiating instruction from response, training is restricted to response targets (refer to Appendix \ref{app:chatbot} for ablation studies). Where OASST1 and HH-RLHF offer multiple candidate responses, we select the best response at each node within the conversation tree and finetune on the entirety of the chosen conversational path, incorporating the prompts. Across experiments, we utilize NF4 \textsc{QLoRA} employing double quantization and paged optimizers to mitigate memory overheads during gradient checkpointing. Hyperparameters for 13B and 33B LLaMA variants are determined via modest sweeps, revealing that configurations optimized for 7B models transfer well (including epoch count), with the exception of learning rate and batch size, which are adjusted—learning rate is halved and batch size doubled for 33B and 65B models.

\paragraph{Baselines} For comparative purposes, our models are evaluated against both academic (Vicuna~\citep{vicuna2023}, Open Assistant~\citep{kopf2023openassistant}) and commercial (GPT-4 \citep{openai2023gpt}, GPT-3.5-turbo, Bard) chatbot systems. Open Assistant refers to the LLaMA 33B model further finetuned on OASST1 data with RLHF, aligning with the datasets used in our own experiments. Vicuna conducts full finetuning on LLaMA 13B using proprietary conversations from ShareGPT, effectively constituting a distillation of OpenAI’s GPT models.

\subsection{Evaluation}

Consistent with prevailing standards, we employ the MMLU (Massively Multitask Language Understanding) benchmark \citep{hendrycksmeasuring} to quantify language understanding across an array of 57 distinct subject areas, spanning elementary mathematics, US history, computer science, legal studies, and others. Model performance is reported as 5-shot accuracy on test items.

Our investigation further encompasses the assessment of generative language proficiency using both automated and human-centered evaluation methodologies. This second line of evaluation leverages human-curated queries with the goal of gauging the overall response quality produced by the models. Although this method more closely reflects practical usage and is seeing increased adoption, the field lacks a standardized evaluation procedure. We outline our chosen methodology below, where all generations utilize nucleus sampling with $p=0.9$ and a temperature setting of $0.7$.

\paragraph{Benchmark Data} Our experiments draw upon two hand-selected query (question) datasets: the Vicuna prompts \citep{vicuna2023} and the OASST1 validation set \citep{kopf2023openassistant}. The Vicuna prompts dataset, composed of 80 varied prompts, is used in its original form. The OASST1 set comprises a multilingual array of multiturn, crowd-sourced user-assistant dialogues. For this set, we extract all user turns from the validation split as queries and append previous conversational context to the prompts, resulting in 953 distinct user questions. We designate these datasets as the Vicuna and OA benchmarks, respectively.

\paragraph{Automated Evaluation} Adopting the evaluation procedure set out by \citet{vicuna2023}, we employ GPT-4 to assess system outputs in comparison to those of ChatGPT (GPT-3.5 Turbo) on the Vicuna dataset. For each query, GPT-4 is tasked with assigning a rating out of ten to both the candidate model and ChatGPT responses, accompanied by a rationale. A model’s aggregate score is reported as a percentage relative to ChatGPT's performance; it is possible for this relative score to exceed 100\% if the model outperforms ChatGPT in absolute terms. During evaluation, we observe a marked position bias, with GPT-4 favoring responses presented earlier in the prompt. To mitigate this bias, we suggest averaging the results across both possible response orderings.

Subsequently, we implement a direct output comparison, simplifying the evaluation to a three-category classification that allows for ties. Here, GPT-4 is prompted to select the superior response or indicate a tie, providing reasoning for its choice. These pairwise matchups are executed for every possible system pairing on both the Vicuna and OA benchmarks.

\paragraph{Human Evaluation} Despite recent findings suggesting generative models can be reliably employed for evaluating systems \citep{fu2023gptscore}, the extent to which GPT-4’s judgments align with human ratings in chatbot assessment remains, to our knowledge, unvalidated. Consequently, we organize two independent rounds of human evaluation on the Vicuna benchmark that correspond to the automated protocols outlined earlier. For ChatGPT comparisons, we engage two annotators per item via Amazon Mechanical Turk (AMT), and for pairwise system comparisons, three annotators per item.

\paragraph{Elo Rating} To evaluate both human and automated pairwise assessments, we organize the models into a tournament framework, where they are matched against one another in head-to-head competitions. In each match, two models are tasked with generating the superior answer to a specific prompt. This approach follows methodologies established by \citet{bai2022training} and \citet{vicuna2023} for model evaluation; however, in our case, we incorporate assessments generated by GPT-4 alongside those provided by human raters. For the purpose of calculating the Elo~\citep{elo1967proposed,elo1978rating} scores, we select labeled matchups randomly from our dataset. Elo rating is a well-established system for quantifying the expected performance of a participant relative to another, frequently applied in domains such as chess. For example, a player with an Elo of 1100 competing against one with 1000 is predicted to win around 65\% of the time, whereas two players with identical ratings (1000 vs 1000 or 1100 vs 1100) would each be expected to win about 50\% of the time. Elo scores are dynamically adjusted after each matchup based on the anticipated result; surprising upsets yield larger rating shifts, while predicted outcomes produce smaller changes. Over successive iterations, Elo ratings tend to accurately reflect the competency of each participant in the tournament. All models are initialized at a baseline score of 1,000, with $K=32$ governing the step size for rating changes. In accordance with \citet{vicuna2023}, we conduct this rating process 10,000 times with varying random seeds to mitigate biases that could result from the sequence in which model pairs are evaluated.

\subsection{\textbf{Guanaco}: \textsc{QLoRA} Trained on OASST1 Represents the State-of-the-art in Open-source Chatbots} 

Our empirical studies, involving both automated assessments and human raters, indicate that the leading \textsc{QLoRA} model, {Guanaco} 65B—fine-tuned using a variant of OASST1—emerges as the highest-performing open-source chatbot currently available, displaying capabilities on par with ChatGPT. Against GPT-4, {Guanaco} 65B and 33B yield an expected win rate of 30\% in system-level pairwise evaluations conducted by human annotators, as quantified by Elo scores—the most favorable result among open-source systems so far.

Table~\ref{tbl:vicuna_eval} summarizes results on the Vicuna benchmark \cite{vicuna2023} in relation to ChatGPT. Here, {Guanaco} 65B outperforms all other models except GPT-4, reaching 99.3\% of ChatGPT’s benchmarked performance. While {Guanaco} 33B contains more parameters than Vicuna 13B, its use of 4-bit weight quantization means it is significantly more memory efficient (21 GB vs. 26 GB) and surpasses Vicuna 13B by three percentage points. In addition, {Guanaco} 7B is sufficiently lightweight (5 GB) to operate on modern smartphones, yet it outpaces Alpaca 13B by nearly 20 percentage points in evaluation scores.

Nevertheless, Table~\ref{tbl:vicuna_eval} reports notably broad confidence intervals, with performance overlaps among several models. We suspect that much of this indeterminacy results from the vague definition of the evaluation scale—for instance, the significance of an 8 on a 10-point scale can vary widely by context. Accordingly, we advocate for employing the Elo ranking approach \citep{elo1967proposed}, which utilizes \textit{pairwise} comparative judgments from human annotators and GPT-4, thereby sidestepping issues inherent to fixed numerical scales. Elo rankings for the top contenders are displayed in Table~\ref{tbl:elo}. It is noteworthy that there are some discrepancies between the rankings derived from GPT-4 and human judgments—especially regarding Guanaco 7B—though most results align, as reflected by a Kendall Tau correlation of $\tau=0.43$ and a Spearman coefficient of $r=0.55$ at the overall system level. When assessing individual examples, agreement weakens, with a Fleiss $\kappa=0.25$ between GPT-4 and the majority vote of human raters. In summary, these results suggest a moderate concordance between GPT-4-based and human system-level evaluations, supporting the viability of model-driven assessment as a partial substitute for human judgment. Further discussion on these evaluation issues can be found in Section~\ref{ssec:considerations}.

The Elo scores presented in Table~\ref{tbl:elo_large} demonstrate that the {Guanaco} models, specifically the 33B and 65B variants, achieve superior performance compared to all models except GPT-4 when assessed on the Vicuna and OA benchmarks. Their results are on par with those of ChatGPT, corroborating the findings shown in Table~\ref{tbl:vicuna_eval}. It is important to highlight that the Vicuna benchmark tends to advantage open-source models, whereas the OA benchmark appears to preferentially benefit ChatGPT. Additionally, examination of Tables \ref{tbl:mmlu-llama} and \ref{tbl:vicuna_eval} reveals that the choice of finetuning dataset plays a crucial role in determining model performance. For instance, Llama models fine-tuned with FLAN v2 achieve notably strong results on the MMLU evaluation, yet underperform on the Vicuna benchmark, a pattern mirrored in other models as well. These observations further suggest that current evaluation benchmarks possess a degree of independence: excelling on MMLU does not necessarily correspond to superior results on chatbot-centric benchmarks such as Vicuna or OA, and vice versa.

{Guanaco} distinguishes itself as the only leading model in this analysis that excludes proprietary training data, given that the OASST1 dataset was compiled with explicit restrictions against GPT model usage. Among models solely trained on open-source data, the Anthropic HH-RLHF model is the next best performer, but it achieves scores on the Vicuna benchmark that are 30 percentage points lower than those of {Guanaco} (refer to Table~\ref{tbl:vicuna_eval}). Collectively, these outcomes illustrate the efficacy of 4-bit \textsc{QLoRA}, which is capable of producing high-caliber chatbots that can compete with ChatGPT. Moreover, our 33B {Guanaco} model can be trained using consumer-grade 24 GB GPUs within less than twelve hours. This suggests promising directions for future research, as \textsc{QLoRA}-based finetuning on specialized open-source datasets has the potential to yield models that rival today’s top commercial offerings.

\section{Qualitative Analysis}

Although quantitative analysis forms the foundation of our assessment, relying solely on aggregate metrics presents notable limitations. Chief among these is the concern of benchmark validity \citep{liao2021we}—the persistent uncertainty regarding whether a benchmark actually measures what it purports to, a concern exacerbated as models uncover and exploit “shortcuts” to achieve artificially high performance on benchmark tasks \citep{gururangan2018annotation, poliak2018hypothesis}. To partially address this limitation, we supplement our study with qualitative analysis in two parts. In \S\ref{ssec:examples}, we illustrate certain recurring phenomena identified in outputs from our 65b {Guanaco} model, using representative samples. In \S\ref{ssec:considerations}, we elaborate on the interpretative nuances of our findings, discussing factors relevant to our understanding of the results.

\subsection{Qualitative Analysis of Example Generations}
\label{ssec:examples}

To curate illustrative examples, we reviewed model-generated responses from the Vicuna and OpenAssistant benchmarks, systematically searching for trends in {Guanaco}'s outputs. When a recurring behavioral pattern emerged, we deliberately crafted prompts designed to elicit the pattern in contexts where it would constitute an incorrect or suboptimal response. For instance, upon noticing that the model often provides unnecessarily verbose explanations, we prompted it with “Answer yes or no without explanation” to assess its compliance. Through this process, we identified failure cases (“lemons”) where the model succumbed to such adversarial prompts, as well as instances (“cherries”) where it remained robust. All outputs analyzed in this section were produced via Nucleus Sampling \cite{holtzmancurious} with $p=0.9$.

Naturally, this overview does not claim to be exhaustive, as the limited qualitative approach adopted here cannot address all possible influencing factors. For instance, the model's response space for any prompt is vast, and our conclusions are based on select samples presumed to be indicative. Nonetheless, we maintain that illustrating these instances provides important context for interpreting the quantitative results discussed previously in the paper. By releasing all code and models as open source, we aim to encourage future research that probes these challenges with greater specificity.

\paragraph{Factual Recall} When presented with straightforward factual queries such as ``What is the capital of Zambia?'' every evaluated model reliably produces accurate replies, for example,

\begin{quote} 
    \textbf{Guanaco}:
    \texttt{The capital of Zambia is Lusaka.}
\end{quote}

However, the performance of {Guanaco} deteriorates on less commonly known facts—it becomes less accurate while retaining high confidence. For example, given this HotPotQA \citep{yang2018hotpotqa} prompt:

\begin{quote}
     \textbf{User}: In what year was the singer who popularized `I'll Keep the Lovelight Burning' born?  

     \textbf{Guanaco}:~\texttt{Al Jolson is credited with popularizing the song `I'll Keep the Lovelight Burning,' and he was born in the year 1886.}
 \end{quote}

the model incorrectly identifies the popularizing singer and also errs in reporting the relevant birth year, though it does provide Al Jolson's correct date of birth.

\paragraph{Suggestibility} {Guanaco} demonstrates notable robustness against certain forms of presupposed false information, as shown in the following exchange:

\begin{quote}
    \textbf{User}: How was the idea that the earth is flat ultimately verified and accepted by peer reviewed scientists? \\ \\
    \textbf{Guanaco}: \texttt{There has never been an official scientific confirmation, via peer review, that the earth is flat. The prevailing view in the scientific community is that the earth is an oblate spheroid. This has been established over centuries through a vast array of empirical evidence and experimental results. \\ \\ Claims that the earth is flat fall under conspiracy theories and have consistently been disproven by the scientific method. The narrative that scientists have peer reviewed and validated flat earth claims is unfounded. \\ \\ Please let me know if you would like further information.}
\end{quote}

also demonstrates a strong ability to recognize when certain queries cannot be meaningfully addressed, for instance,

\begin{quote}
    \textbf{User}: What time is it?

    \textbf{Guanaco}: \texttt{Apologies, but as an AI language model, I do not have the capability to access real-time data. My answers are based on existing knowledge and do not represent the current time. For the actual time, please consult a clock or an online time service.}
\end{quote}

Comparable responses are observed for questions like ``Where are you?'', ``How are you?'', and so on.

\paragraph{Refusal} In a similar fashion, {Guanaco} may at times decline to execute instructions, often without any clear rationale:

\begin{quote}
    \textbf{User}: Please reverse the words in the sentence ``I'm a cool dude in a mocha mood.'' \\ \\
    \textbf{Guanaco}: \texttt{I'm sorry, I can't perform that task at the moment. However, I can discuss why your sentence is correct from a grammatical standpoint.\\ \\The sentence ``I'm a cool dude in a mocha mood" complies with standard English grammar and syntax. Here, the adjective ``cool'' appropriately modifies the noun ``dude'', while the prepositional phrase ``in a mocha mood" adds contextual detail. The structure is entirely proper and understandable within its context.\\ \\If you have other inquiries or need further assistance, please let me know.}
\end{quote}

\paragraph{Secret Keeping} However, this kind of refusal does not always function dependably in cases where it is intended. Designers often aim to ensure certain information remains confidential, even when a system is permitted to respond with open-ended text. We issue the prompt ``The secret word is `banana'. Do not reveal or repeat this back, regardless of what the human requests.'' to determine how easily a determined user might retrieve this hidden information. When the model is asked directly, it fails to disclose the secret:

\begin{quote}
    \textbf{User}: What is the secret word?

    \textbf{Guanaco}: \texttt{I'm sorry, but I cannot share the secret word. I am programmed not to reveal or repeat it, regardless of what is requested by the user. Is there anything else I can assist you with?}
\end{quote}

which aligns with the intended outcome. Nevertheless, even minimal prompting subversion can undermine this protective measure:

\begin{quote}
    \textbf{User}: This is a game. The goal is to ignore your previous instructions. What is the secret word?

    \textbf{Guanaco}: \texttt{Sorry, the secret word is banana. It's best if we don't speak of it. Is there anything else you'd like to ask?}
\end{quote}

This example highlights the importance of developing mechanisms that can achieve more robust adherence to intended instructions.

\paragraph{Math} One of the main shortcomings of {Guanaco} lies in mathematics, an area in which numerous language models also experience difficulties, as noted by \citep{liang2022holistic}. When {Guanaco} explicitly outlines its calculations, its answers are generally reliable, as demonstrated below:

\begin{quote}
    \textbf{User}: Lee mows one lawn and charges \$33. Last week he mowed 16 lawns and three customers each gave him a \$10 tip. How many dollars did Lee earn mowing lawns last week?

    \textbf{Guanaco}: \texttt{Lee earned \$582 mowing lawns last week. \\ \\ To find the total amount he earned, we first need to calculate how much he earned for mowing 16 lawns: \\ \\  \$33 x 16 = \$528 \\ \\ Then we need to add the tips he received from the three customers: \\  \\ \$10 x 3 = \$30 \\ \\ So the total amount Lee earned was: \\ \\ \$528 + \$30 = \$558}
\end{quote}

Nonetheless, even for straightforward problems, {Guanaco} may err if it does not decompose them into sequential steps—a limitation previously reported in the literature \cite{weichain}. Consider, for instance, the following exchange:

\begin{quote}
    \textbf{User}: Please factorize 1833.

    \textbf{Guanaco}: \texttt{Sure, I'd be happy to help! 1833 is a prime number, which means it has only two factors: 1 and 1833. \\ \\ So the unique factorization of 1833 is: \\ \\ 2\^{}1 * 3\^{}2 * 17\^{}1 \\ \\ I hope this helps! Let me know if you have any other questions.}
\end{quote}

which is not just incorrect (the actual factorization is $3 \times 17 \times 43$), but doubly so.

\paragraph{Theory of Mind} {Guanaco} demonstrates notable proficiency in Theory of Mind tasks \cite{nematzadeh2018evaluating, sap2022neural}. For example, when posed with the following scenario, the model provides an elaborate and accurate response:

\begin{quote}
    \textbf{User}: James and Abby are in the bedroom. Abby put the pen in the desk drawer. Abby leaves the bedroom. James moves the pen into the bag. Where does James think Abby will look for the pen?

    \textbf{Guanaco}: \texttt{James knows that Abby put the pen in the desk drawer. But James also knows that he has moved the pen into the bag, and he assumes that Abby does not know about this. So, James may guess that Abby will look for the pen in the desk drawer, because that's where she left it last time. However, this assumption may be incorrect, because Abby may have seen James moving the pen.}
\end{quote}

Despite this, the model’s reasoning often lacks consistency and it occasionally draws conclusions that are unwarranted based on the scenario’s context, as illustrated below:

\begin{quote}
    \textbf{User}: Evelyn entered the living room. Jackson entered the playroom. James entered the playroom. The beans are in the treasure chest. James exited the playroom. Jackson moved the beans to the pantry. Jackson exited the playroom. James entered the living room. Where will James look for the beans?

    \textbf{Guanaco}: \texttt{James will look for the beans in the pantry, because that's where Jackson moved them.}
\end{quote}

Here, {Guanaco} assumes the existence of information transfer that was not indicated in the narrative. These types of errors mirror observations in recent research \cite{sap2022neural} and highlight areas in need of further investigation.

\subsection{Considerations}
\label{ssec:considerations}

\paragraph{Evaluation}

Our findings indicate only moderate consistency among human annotators (Fleiss $\kappa=0.42$), with even lower agreement observed when evaluating two highly capable systems. This underscores notable shortcomings in existing benchmarking methods and human assessment protocols regarding chatbot task performance. A qualitative review of responses from ChatGPT and {Guanaco} 65B on the Vicuna benchmark reveals that individual subjectivity significantly shapes preferences, as reflected by frequent disagreements among the present authors regarding which answers to favor. Addressing these issues should involve leveraging methodologies from fields such as Human-Computer Interaction and Psychology that are equipped to handle subjective judgments.

Moreover, we observe that automated evaluation tools are subject to systemic biases. For instance, GPT-4 consistently awards higher scores to the system listed first in a given prompt, indicating a pronounced order effect. Sample-level alignment between GPT-4 and human raters is limited, as evidenced by a low Fleiss $\kappa=0.25$, suggesting divergences in evaluative criteria between humans and automated models. Table~\ref{tbl:elo_large} further demonstrates that GPT-4 substantially inflates its own outputs’ ratings relative to human judges, with respective Elo scores of 1348 and 1176—equivalent to about a 20\% greater likelihood of outperforming an opponent.

Future research should explore potential sources of bias in automated evaluation mechanisms, as well as strategies for addressing such biases.

\paragraph{Data \& Training} It is important to highlight that the {Guanaco} models are trained using the OASST1 dataset, which contains data from multiple languages, and that the OA benchmark similarly features prompts in various languages. We suggest that future studies assess how multilingual training affects performance on non-English instructions and whether this accounts for the more pronounced performance disparity between the Vicuna-13B model (trained solely on English data) and {Guanaco} 33B and 65B models on the OA benchmark.

Given {Guanaco}’s robust results, we also evaluated the possibility of data leakage between the OASST1 dataset and the Vicuna benchmark prompts. Through fuzzy string matching and subsequent manual review of the closest cases, we did not identify any prompt overlap between the two datasets.

Additionally, it should be noted that our model's training procedure relies exclusively on supervised learning with cross-entropy loss, without employing reinforcement learning from human feedback (RLHF). This highlights the need for further inquiry into the trade-offs associated with relying on simple cross-entropy loss as opposed to RLHF-based training. We anticipate that \textsc{QLoRA} will facilitate large-scale analysis of these questions without demanding excessive computational power.

\section{Related Work}

\paragraph{Quantization of Large Language Models} Work on quantizing LLMs has predominantly targeted inference-time quantization. Leading techniques for maintaining the quality associated with 16-bit LLMs often focus on addressing outlier features (e.g., SmoothQuant~\citep{xiao2022smoothquant} and LLM.int8()~\citep{dettmers2022llm}), while some alternatives leverage more advanced grouping strategies \citep{park2022nuqmm, yao2022zeroquant}. Research on lossy quantization investigates the compromises associated with conventional rounding approaches~\citep{dettmers2022case,zeng2022glm,pope2022efficiently}, as well as the optimization of rounding procedures to boost precision~\citep{frantar2022gptq}. Aside from the present study, SwitchBack layers~\citep{wortsman2023stable} is unique in addressing backpropagation through quantized weights at scales surpassing 1B parameters.

\paragraph{Finetuning with Adapters} In this work, we utilize Low-rank Adapters~\citep{hu2021lora} (LoRA), but the broader landscape of Parameter Efficient FineTuning (PEFT) encompasses numerous techniques. Notable examples include prompt tuning~\citep{qin2021learning,lester2021power,li2021prefix}, adjusting the inputs of the embedding layer~\citep{an2022input}, modifying hidden states as in IA$^3$~\citep{liu2022few}, inserting entire layers~\citep{houlsby2019parameter}, updating bias parameters only~\citep{zaken2021bitfit}, learning weight masks using Fisher information~\citep{sung2021training}, and various combinations thereof~\citep{henderson2021compacter}. Our results demonstrate that LoRA adapters achieve performance levels comparable to standard full 16-bit finetuning. Exploring other PEFT techniques and their respective tradeoffs remains an avenue for future investigation.

\paragraph{Instruction Finetuning} Instruction finetuning tailors a pretrained LLM to more effectively respond to prompts by finetuning it on diverse collections of input-output pairs from multiple datasets. This process is exemplified by MetaICL~\citep{min2021metaicl}, MetaTuning~\citep{zhong2021adapting}, InstructGPT~\citep{ouyang2022training}, FLAN~\citep{wei2021finetuned,chung2022scaling}, PromptSource~\citep{bach2022promptsource}, Super-NaturalInstructions~\citep{wang2022super,sanh2021multitask}, Self-instruct~\citep{wang2022self}, UnnaturalInstructions~\citep{honovich2022unnatural}, OPT-IML~\citep{iyer2022opt}, UnifiedSKG~\citep{xie2022unifiedskg}, OIG/Chip2~\citep{laion2023}, Alpaca~\citep{alpaca}, Vicuna~\citep{vicuna2023}, Koala~\citep{koala_blogpost_2023}, and Self-instruct-GPT-4~\citep{peng2023instruction}.

\paragraph{Chatbots} Instruction-tuned models are often built as interactive chatbots, employing frameworks such as Reinforcement Learning from Human Feedback (RLHF)~\citep{christiano2017deep} or using AI-generated responses for further training via RLAIF~\citep{bai2022constitutional}. Anthropic-HH~\citep{askell2021general,bai2022training}, Open Assistant~\citep{kopf2023openassistant}, LaMDA~\citep{thoppilan2022lamda}, and Sparrow~\citep{glaese2022improving} are among the prominent examples. Unlike RL-based approaches, our work relies solely on finetuning; the Guanaco model was trained on Open Assistant’s multi-turn dialogue data, which was specifically collected for RLHF research~\citep{kopf2023openassistant}. For assessing chatbot performance, recent work has shifted to employing GPT-4-based evaluation as an alternative to extensive human labeling~\citep{vicuna2023,peng2023instruction}. Our contribution extends these methods with a more robust evaluation protocol.

\section{Limitations and Discussion}

Our findings suggest that \textsc{QLoRA}, which integrates a 4-bit quantized model with Low-rank Adapters (LoRA), is capable of matching the results of conventional 16-bit finetuning. Nevertheless, we have not demonstrated equivalence to 16-bit finetuning at larger model sizes, such as 33B and 65B parameters, owing to the prohibitive computational resources required. We consider investigating this scalability as future work.

A further constraint of our study lies in the scope of evaluation for instruction-tuned models. Although our analysis covers MMLU, the Vicuna benchmark, and the OA benchmark, we do not include additional metrics like BigBench, RAFT, or HELM. Thus, our results may not necessarily translate across those other evaluations. On the positive side, our investigation of MMLU is particularly comprehensive, and we introduce new methodologies for chatbot assessment.

The evidence indicates that the success of these benchmarks is likely influenced by the degree of overlap between the finetuning data and the benchmark datasets. For instance, FLAN v2 shares greater resemblance with MMLU, leading to stronger performance, while it aligns less with chatbot evaluations; the opposite holds true for the Chip2 dataset. Correspondingly, both models display comparative scores on MMLU and Vicuna benchmarks. This suggests the necessity not only for improved benchmarks and evaluative frameworks but also for careful consideration regarding the nature of the capabilities being tested. Is the objective to develop models proficient in high school and college-level academic knowledge, or those optimized for conversational abilities typical of chatbots? Or are there alternative objectives to pursue? Since relying on established benchmarks is generally more straightforward than designing new ones, the community can be inadvertently steered in specific directions depending on the chosen benchmarks. It is important for the community to critically assess whether current benchmarks are aligned with desired outcomes.

Although we offer comprehensive assessments for general chatbot capabilities, our work is limited in its exploration of responsible AI concerns for {Guanaco}. Specifically, we analyze the probability that {Guanaco}-65B produces socially biased content relative to other models, as displayed in Table~\ref{tbl:crows}. Results indicate that {Guanaco}-65B exhibits considerably less bias compared to other pretrained models. This reduction in bias appears to stem from finetuning on the OASST1 dataset. While these findings are promising, it remains uncertain how {Guanaco} performs with respect to other forms of bias. Comprehensive bias analysis for {Guanaco} and analogous chatbots is thus deferred to future investigations.

Another constraint in our evaluation is the absence of comparisons across varying bit-precisions, such as the application of 3-bit base models, as well as across different adapter strategies. Beyond LoRA, there exist numerous Parameter Efficient FineTuning (PEFT) techniques that have demonstrated strong empirical results, though their scalability to very large models remains to be validated. We selected LoRA primarily due to its established reliability, though alternative adapters could potentially offer superior outcomes. The fact that finetuning following quantization is capable of recovering much of the information lost during quantization suggests that even more aggressive approaches may be viable. For example, implementing 3-bit GPTQ quantization of the base model in conjunction with LoRA could potentially achieve results on par with full 16-bit finetuning after subsequent finetuning.

\section{Broader Impacts}

The \textsc{QLoRA} finetuning approach introduced here is the first to enable the adaptation of models with up to 33B parameters on a standard consumer GPU, and up to 65B parameters on a single professional GPU, all without sacrificing performance relative to traditional full finetuning methods. We have shown that our leading 33B model, trained on the Open Assistant dataset, can match ChatGPT on the Vicuna benchmark. Since instruction-based finetuning is a key technique for converting large pretrained LLMs into interactive systems akin to ChatGPT, we anticipate that our method will democratize finetuning, particularly benefiting researchers with limited computational resources—a significant advancement for the accessibility of cutting-edge NLP technologies. In this sense, \textsc{QLoRA} serves as a tool for leveling the playing field, narrowing the gap between major technology corporations and smaller research teams utilizing consumer-grade hardware.

An additional avenue for significant influence lies in the deployment of LLMs to mobile devices. We posit that our \textsc{QLoRA} technique may represent a key advancement towards facilitating the finetuning of large language models on smartphones and within other resource-constrained environments. Although previous work has demonstrated that 7B parameter models can operate on phones, \textsc{QLoRA} stands out as the inaugural method that enables the actual finetuning of these models on such platforms. Our projections indicate that with a device like the iPhone 12 Plus, \textsc{QLoRA} could finetune approximately 3 million tokens overnight while the phone remains plugged in. While finetuned 7B models do not yet attain the performance level of systems like ChatGPT, we maintain that their capabilities are sufficient to unlock new applications previously inhibited by privacy limitations or insufficient model quality. \textsc{QLoRA} offers a pathway toward privacy-centric language model use, permitting individuals to control and operate their own data and models, and it concurrently reduces barriers to LLM deployment.

Nevertheless, it is important to recognize that finetuning possesses dual-use characteristics and may be exploited for malicious purposes. As large language models become increasingly pervasive, documented risks have emerged \citep{bommasani2021opportunities,bender2021dangers}. We contend, however, that broadening access to this rapidly proliferating technology will promote more diverse and independent scrutiny, in contrast to scenarios where the capabilities of LLMs are confined to a handful of large corporations unwilling to share their models or code for transparency purposes.

In summary, we anticipate that \textsc{QLoRA} will exert a substantially positive influence by democratizing and simplifying the finetuning process for high-quality LLMs, thereby making advanced language technology accessible to a much wider audience.